\documentclass[letterpaper,10pt,conference]{ieeeconf}

\IEEEoverridecommandlockouts
\overrideIEEEmargins

% ============================================================
% STEER -- ICRA 2027 DRAFT
%
% Working title:
% STEER: Cross-Backbone Causal Temporal Refinement
% for Frozen Surgical Stereo
%
% IMPORTANT:
% - Canonical model: seed-0, bounded H=4 operating policy.
% - H=4 is a frozen risk/gain operating point, NOT the minimum-EPE
%   horizon on D2.
% - D2 is development.
% - D7 is held-out confirmation.
% - DRENDS zero-shot results use the SCARED-trained canonical model.
% - Runtime/downstream/clinical claims remain unsupported until tested.
% ============================================================

\usepackage[T1]{fontenc}
\usepackage[utf8]{inputenc}

\usepackage{amsmath}
\usepackage{amssymb}
\usepackage{bm}

\usepackage{graphicx}
\usepackage{booktabs}
\usepackage{multirow}
\usepackage{array}
\usepackage{siunitx}

\usepackage{algorithm}
\usepackage{algpseudocode}

\usepackage{xcolor}
\usepackage{cite}
\usepackage{url}
\usepackage{balance}

% ============================================================
% DRAFT HELPERS
% ============================================================

\newif\ifinternaldraft
\internaldrafttrue
% \internaldraftfalse

\newcommand{\TODO}[1]{%
    \ifinternaldraft
        \textcolor{red}{\textbf{[TODO: #1]}}%
    \fi
}

\newcommand{\NOTE}[1]{%
    \ifinternaldraft
        \textcolor{blue}{\textbf{[NOTE: #1]}}%
    \fi
}

% ============================================================
% RED DRAFT MARKERS
% ============================================================
% All \TODO{...} blocks are intentionally rendered in red while
% \internaldrafttrue. They denote experiments/results that are not yet
% available and MUST NOT be converted into prose claims until verified.
%
% ============================================================
% METHOD / METRIC MACROS
% ============================================================

\newcommand{\steer}{\textsc{STEER}}
\newcommand{\raw}{\mathrm{raw}}
\newcommand{\mem}{\mathrm{mem}}
\newcommand{\fused}{\mathrm{fused}}
\newcommand{\warp}{\mathcal{W}}

\newcommand{\EPE}{\mathrm{EPE}}
\newcommand{\HUR}{\mathrm{HUR}}
\newcommand{\BUR}{\mathrm{BUR}}
\newcommand{\HPlus}{H^{+}}
\newcommand{\BPlus}{B^{+}}
\newcommand{\Bad}{\mathrm{Bad}}
\newcommand{\DTCE}{\mathrm{DTCE}}
\newcommand{\TEPE}{\mathrm{TEPE}}
\newcommand{\OPW}{\mathrm{OPW}}
\newcommand{\AURC}{\mathrm{AURC}}

\newcommand{\R}{\mathbb{R}}

% ============================================================
% TITLE
% ============================================================

% \NOTE: NAME COLLISION. "STEER" is already used in robotics by
% "STEER: Flexible Robotic Manipulation via Dense Language Grounding",
% arXiv:2411.03409 (Nov 2024), which targets the same venue population.
% No collision was found in the stereo/depth literature. The acronym is
% also never expanded anywhere in this paper. Rename before submission,
% or expand it and cite the collision explicitly.

\title{
\textbf{STEER: Cross-Backbone Causal Temporal Refinement\\
for Frozen Surgical Stereo}
}

\author{Anonymous Authors}

\begin{document}

\maketitle
\thispagestyle{empty}
\pagestyle{empty}

% ============================================================
% ABSTRACT
% ============================================================

\begin{abstract}

Dense stereo is attractive for robotic surgery because it provides
metric geometry without active illumination, yet modern stereo
estimators are typically applied independently to each frame and can
therefore fluctuate under specularities, weak texture, tissue
deformation, tool occlusions, and camera motion.
We ask whether temporal information can improve frozen stereo
predictions without modifying the stereo estimator, accessing its
internal representation, or using future frames.

We present \steer, a causal plug-and-play temporal refiner for frozen
stereo.
At each frame, preceding disparity is motion-aligned to the current view
using frozen SEA-RAFT.
A shared 142-channel external evidence representation feeds a
177,338-parameter CODD-inspired reset-and-fusion head that predicts a
bounded interpolation between the current stereo estimate and aligned
temporal memory.
Corrected recurrence is explicitly bounded and periodically re-anchored
to independently predicted raw geometry.

A controlled development-set closure shows that the learned policy is
not explained by temporal smoothing:
on SCARED-C D2, canonical $H=4$ reduces mean EPE by $4.05\%$, compared
with only $0.59\%$ for the best matched fixed/EMA blend, while direct
flow-confidence weighting and ungated memory propagation degrade
accuracy.
The raw-versus-memory oracle exposes a $12.30\%$ improvement ceiling,
of which STEER recovers $32.9\%$.
Continuous corrected-state recurrence increases EPE by $62.69\%$,
supporting bounded memory.

One frozen $H=4$ checkpoint subsequently improves all five evaluated
stereo estimators on held-out SCARED-C D7, including CREStereo and
Fast-FoundationStereo, which were excluded from temporal-head training.
Mean backbone EPE decreases from $0.4921$ to $0.4696$, while Bad1 and
disparity RMSE also decrease for every backbone.
Zero-shot evaluation on five DRENDS recordings ($20{,}929$ frames) with
three estimators, two of them unseen in training, improves disparity
EPE, Bad3 and metric depth in all $60$ recording--backbone--metric
cells, completing the transfer matrix across both the backbone and the
domain shift.
We further show that confining the fused output to the support the
module already declares is load-bearing rather than cosmetic: without
it, zero-padded out-of-bounds memory turns rare pixels into invalid
predictions and produces the metric-depth tail previously attributed to
the external domain.

We additionally introduce a unified refinement-centric evaluation
framework that jointly measures spatial geometry, multi-horizon
temporal fidelity, introduced and recovered error, frame-level tails,
and selective risk under fixed prediction-independent support.

\end{abstract}


% ============================================================
% I. INTRODUCTION
% ============================================================

\section{Introduction}
\label{sec:intro}

Metric three-dimensional perception is a fundamental component of
robot-assisted surgery.
Dense geometry can support camera navigation, tool--tissue distance
estimation, augmented reality, deformation analysis, and autonomous
manipulation.
Stereo endoscopy is particularly attractive because depth can be
recovered passively from the images already available on many surgical
robot platforms.

Surgical stereo nevertheless remains difficult.
Endoscopic scenes contain weak or repetitive texture, non-Lambertian
tissue, strong specular highlights, blood, smoke, deformable anatomy,
rapid viewpoint changes, and frequent tool occlusions
\cite{allan2021scared,edwards2022servct}.
Modern stereo estimators such as RAFT-Stereo
\cite{lipson2021raftstereo}, CREStereo
\cite{li2022crestereo}, Stereo Anywhere
\cite{bartolomei2025stereoanywhere}, S$^2$M$^2$
\cite{min2025s2m2}, and recent foundation models
\cite{wen2025foundationstereo,wen2026fastfoundationstereo}
provide increasingly strong frame-wise disparity.
However, a frame-wise model has no explicit requirement that predictions
remain coherent over time.

Video stereo methods address this limitation by integrating temporal
information directly into the model.
DynamicStereo aggregates spatio-temporal information over stereo video
\cite{karaev2023dynamicstereo}.
CODD combines stereo estimation, learned motion, and temporal fusion in
an online framework \cite{li2023codd}.
Match Stereo Videos and its plug-in BiDAStabilizer employ bidirectional
alignment and sequence propagation \cite{jing2024bida}.
PPMStereo uses learned memory construction for long-range dynamic stereo
\cite{wang2025ppm}.
These methods establish the value of temporal information, but commonly
assume a particular temporal architecture, model-specific
representations, future observations, or co-design with the stereo
estimator.

We study a different deployment setting:

\begin{quote}
\emph{Can one causal temporal module be trained once and attached
unchanged to heterogeneous frozen stereo estimators?}
\end{quote}

This formulation matters for robotic systems because the stereo
backbone and temporal logic can evolve independently.
A new stereo estimator can replace the previous one without requiring
the temporal module to access its cost volume, hidden state, confidence
head, or training pipeline.

Our starting point is the reset-and-fusion principle introduced by CODD
\cite{li2023codd}.
We do not claim reset/fusion itself as novel.
Instead, we retain its useful inductive bias:
the temporal system should decide whether motion-aligned historical
geometry should influence the current stereo estimate and constrain the
correction to a convex interpolation between the two.

We redesign this principle for an external black-box interface.
The integrated motion pathway is replaced by frozen SEA-RAFT;
stereo-backbone-specific representations are replaced by a shared
external evidence space; and the same learned temporal head is evaluated
unchanged across heterogeneous stereo architectures.

The resulting system has a deliberately constrained structure.
The current stereo estimate remains an anchor.
Temporal memory is geometrically aligned before use.
The learned head controls the amount of intervention, but cannot produce
a value outside the interval spanned by the current and temporal
candidates.
Finally, corrected recurrence is periodically reset to raw stereo
geometry.

These restrictions raise an important question:
does the learned head actually provide value, or would simple temporal
averaging give the same improvement?

We answer this experimentally under a frozen development protocol.
Matched fixed fusion, EMA, warped-memory propagation, and direct
forward--backward confidence weighting recover only a small fraction of
STEER's gain.
In contrast, a ground-truth oracle shows that the aligned memory contains
substantially more useful information than the current policy exploits.
This indicates that the core problem is not obtaining temporal evidence,
but identifying \emph{where and how much} of it should be used.

A second question is how such intervention should be evaluated.
Temporal smoothness is not equivalent to temporal correctness:
a stale prediction can be smooth while being geometrically wrong.
Likewise, lower mean EPE does not reveal whether refinement introduces
new bad pixels or rare catastrophic frames.
We therefore evaluate refinement relative to the frozen upstream
prediction and explicitly separate geometry, temporal fidelity,
benefit, harm, and tail behavior.

Our contributions are:

\begin{itemize}

    \item We formulate surgical temporal stereo refinement as a
    \emph{causal plug-and-play adaptation problem} over frozen,
    heterogeneous stereo estimators, requiring neither future frames
    nor access to backbone-specific internal representations.

    \item We introduce \steer, a 177,338-parameter CODD-inspired
    temporal decision module combining causal SEA-RAFT alignment, a
    shared 142-channel external evidence space, convex disparity
    fusion, and bounded corrected-state recurrence.

    \item Through a controlled matched-baseline closure, we show that
    STEER is not explained by simple temporal smoothing:
    canonical $H=4$ achieves a $4.05\%$ EPE reduction compared with
    $0.59\%$ for the best fixed/EMA blend, while unrestricted recurrence
    fails strongly.
    An oracle analysis reveals substantial remaining temporal headroom.

    \item A single frozen $H=4$ checkpoint improves all five evaluated
    stereo estimators on held-out SCARED-C D7, including two
    architectures excluded from temporal-head training, and shows
    preliminary zero-shot external-domain transfer on DRENDS with a
    training-seen backbone.

    \item We introduce a unified evaluation framework for temporal
    refinement that jointly measures spatial accuracy, multi-horizon
    temporal fidelity, introduced harm and recovered error, frame-level
    failure tails, and optional selective risk under fixed,
    prediction-independent support.

\end{itemize}


% ============================================================
% II. RELATED WORK
% ============================================================

\section{Related Work}
\label{sec:related}

\subsection{Deep Stereo Matching}

Deep stereo matching has progressed from learned matching costs toward
iterative recurrent estimation and broad-domain zero-shot models.
RAFT-Stereo adapts recurrent all-pairs field transforms to rectified
stereo and performs iterative disparity refinement
\cite{lipson2021raftstereo}.
CREStereo combines cascaded recurrent estimation with adaptive
correlation \cite{li2022crestereo}.

Stereo Anywhere combines stereo evidence with monocular foundation-model
priors \cite{bartolomei2025stereoanywhere}, while FoundationStereo and
Fast-FoundationStereo target broad cross-domain generalization
\cite{wen2025foundationstereo,wen2026fastfoundationstereo}.
S$^2$M$^2$ uses scalable multi-resolution matching
\cite{min2025s2m2}.

Our work is orthogonal to these developments.
We neither modify the stereo estimator nor require a particular stereo
architecture.


\subsection{Temporal Stereo and Depth}

CODD estimates temporally consistent online depth by combining stereo,
motion, and fusion modules and introduced a lightweight reset/fusion
mechanism for deciding when historical geometry should influence the
current estimate \cite{li2023codd}.

DynamicStereo jointly processes stereo video using spatio-temporal
attention \cite{karaev2023dynamicstereo}.

Match Stereo Videos introduces bidirectional adjacent-frame alignment
and global bidirectional propagation; BiDAStabilizer exposes temporal
refinement as a plug-in stabilizer for image-based stereo methods
\cite{jing2024bida}.
Its future-frame and backward-sequence processing distinguish the full
formulation from strictly causal robotic deployment.

PPMStereo addresses longer temporal dependencies through learned
pick-and-play memory construction \cite{wang2025ppm}.
Its memory operates inside an integrated dynamic-stereo architecture
rather than through an external interface to arbitrary frozen stereo
estimators.

Plug-and-play stabilization has also been investigated for monocular
video depth.
NVDS and related neural video-depth stabilizers refine outputs from
single-image depth estimators \cite{wang2023nvds}.
Monocular relative-depth stabilization and calibrated stereo disparity,
however, operate in different output spaces.


\subsection{Evaluating Temporal Refinement}

Stereo benchmarks predominantly report frame-wise geometric quantities
such as EPE and bad-pixel rates.
Video methods additionally report temporal-consistency measures.
Neither family alone fully characterizes a temporal intervention.

A prediction can become temporally smoother without becoming more
geometrically correct.
Conversely, a refiner can reduce average geometric error while
introducing new failures on initially correct pixels.

Temporal refinement is inherently comparative:
every refined output is an intervention on an existing upstream
prediction.
This motivates explicit measurement of both recovered and introduced
error.

When a method exposes a confidence or risk score, evaluation further
connects to selective prediction and reject-option learning
\cite{geifman2019selectivenet}.

Our contribution is not that every constituent metric is individually
new.
Instead, we organize established spatial, temporal, risk, and
refinement-relative quantities into a common protocol under fixed
prediction-independent support.


\subsection{Surgical Stereo Geometry}

SCARED provides stereo correspondence and reconstruction data for
endoscopic scenes \cite{allan2021scared}.
SERV-CT provides CT-derived reference disparities and depths for
endoscopic 3D reconstruction validation \cite{edwards2022servct}.
Dynamic surgical reconstruction has also been studied through dedicated
stereo pipelines \cite{long2021edssr}.

D4D contains deformable surgical sequences with complementary geometric
acquisition \cite{docea2026d4d}.
DRENDS targets depth estimation under dynamic robotic endoscopy and
provides stereo imagery, calibration, and metric depth observations
under tissue deformation \cite{loza2025drends}.


% ============================================================
% III. METHOD
% ============================================================

\section{Method}
\label{sec:method}

\subsection{External Causal Interface}

Let frozen stereo estimator $S$ process the current rectified pair:
%
\begin{equation}
d_t^{\raw},M_t^{\raw}
=
S(I_t^L,I_t^R),
\label{eq:stereo}
\end{equation}
%
where $d_t^{\raw}$ is positive-left disparity and $M_t^{\raw}$ its
validity mask.

The temporal module never accesses the internal representation of
$S$.
It consumes only externally observable quantities:
RGB, disparity, validity, optical flow, and derived evidence.

At time $t$, the cache contains the preceding left RGB image, preceding
raw disparity and validity, the previous fused state, and its state age.
No observation from $t+1$ or later is available.

The goal is
%
\begin{equation}
d_t^{\fused}
=
T_\theta
\left(
d_t^{\raw},
m_{t-1},
I_t^L,
I_t^R,
I_{t-1}^L
\right).
\label{eq:adapter}
\end{equation}


\subsection{Causal Motion Alignment}

Frozen SEA-RAFT \cite{wang2024searaft} estimates both geometric flow
directions:
%
\begin{align}
F_{t\rightarrow t-1}
&=
F(I_t^L,I_{t-1}^L),
\\
F_{t-1\rightarrow t}
&=
F(I_{t-1}^L,I_t^L).
\end{align}

The second direction is used only to measure forward--backward
reliability.
Both flows depend exclusively on two already observed images.

Previous geometry is pulled to the current image grid:
%
\begin{equation}
\widetilde m_t(x)
=
\warp
\left(
m_{t-1},
F_{t\rightarrow t-1}
\right)(x)
=
m_{t-1}
\left(
x+F_{t\rightarrow t-1}(x)
\right).
\label{eq:warp}
\end{equation}

Bilinear sampling with zero padding and
\texttt{align\_corners=True} is used.

Warp support $S_t$ indicates whether the source coordinate lies inside
the preceding image.
Historical validity is sampled using the same grid.

The reverse flow is pulled to the current grid:
%
\begin{equation}
\widetilde F_{t-1\rightarrow t}
=
\warp
\left(
F_{t-1\rightarrow t},
F_{t\rightarrow t-1}
\right).
\end{equation}

Forward--backward residual is
%
\begin{equation}
e_t^{FB}
=
\left\|
F_{t\rightarrow t-1}
+
\widetilde F_{t-1\rightarrow t}
\right\|_2.
\end{equation}

The corresponding reliability cue is
%
\begin{equation}
C_t^{FB}
=
S_t
\exp
\left(
-\frac{e_t^{FB}}
{
0.5+
0.01
\left(
\|F_{t\rightarrow t-1}\|_2+
\|\widetilde F_{t-1\rightarrow t}\|_2
\right)
}
\right).
\label{eq:fb}
\end{equation}

Bidirectional \emph{flow estimation} between two already observed
frames does not constitute bidirectional temporal inference.


\subsection{Shared External Evidence}

All stereo estimators are mapped into one common evidence representation:
%
\begin{equation}
C_t
=
\Phi
\left(
d_t^{\raw},
\widetilde m_t,
I_t^L,
I_t^R,
I_{t-1}^L,
F_{t\rightarrow t-1},
C_t^{FB},
S_t,
\widetilde M_t
\right).
\end{equation}

The canonical representation consumed by the decision head contains
142 channels, but the evidence is constructed at two resolutions.
At the one-quarter-resolution evidence grid, 129 channels encode
matching/correlation and frozen appearance information.
These include matching curves around raw and memory disparity,
disparity and appearance self-correlations, cross-correlations, sampled
matching-support cues, and 64 frozen ResNet-18 layer-1 channels
\cite{he2016resnet}.

A separate 13-channel cache-grid group encodes current disparity,
aligned memory, signed and absolute disagreement, optical-flow
components, flow magnitude, forward--backward confidence, warp support,
aligned validity, and RGB evidence.
The 129-channel quarter-resolution group is brought to the cache-grid
resolution before concatenation with these 13 channels; hence the
decision head receives a single 142-channel tensor.

Normalization is fixed before training.
Disparities are clipped to $[0,64]$ and scaled by $64$;
signed disparity residuals are clipped to $[-16,16]$ and scaled by
$16$;
flow components use $[-32,32]$ and flow magnitude $[0,32]$.
Frozen image features use ImageNet-normalized RGB.


\begin{figure*}[t]
\centering

\IfFileExists{figures/steer_overview.pdf}{
\includegraphics[width=0.96\textwidth]{figures/steer_overview.pdf}
}{
\fbox{
\parbox{0.94\textwidth}{
\centering
\small
\vspace{1.5mm}
Frozen stereo $d_t^{raw}$ + bounded state $m_{t-1}$
$\rightarrow$ SEA-RAFT alignment
$\rightarrow \widetilde m_t$
\\[1.5mm]
$\downarrow$
\\[1.5mm]
142-channel external evidence
$\rightarrow$ learned reset $\times$ fusion
$\rightarrow w_t$
\\[1.5mm]
$\downarrow$
\\[1.5mm]
$d_t^{STEER}=(1-w_t)d_t^{raw}+w_t\widetilde m_t$
$\rightarrow$ bounded writeback + periodic raw re-anchor
\vspace{1.5mm}
}}
}

\caption{
STEER uses frozen stereo and frozen motion estimation.
Only the compact reset-and-fusion decision head is trained.
}
\label{fig:overview}

\end{figure*}


\subsection{Reset-and-Fusion Decision Head}

The decision structure is inspired by CODD \cite{li2023codd}.
We do not claim reset/fusion or convex interpolation as novel in
isolation.

The cache-grid pathway maps the 142-channel input to 24 channels through
a $3\times3$ convolution, GroupNorm, SiLU, and a residual block.
A coarse pathway expands to 48 channels at one-quarter resolution,
followed by three 48-channel residual blocks.

A full-resolution reset branch predicts
%
\begin{equation}
r_t=\sigma(z_t^r),
\end{equation}
%
while the coarse fusion branch predicts
%
\begin{equation}
f_t=
\uparrow_4[\sigma(z_t^f)].
\end{equation}

The effective temporal weight is
%
\begin{equation}
w_t=r_tf_t.
\label{eq:weight}
\end{equation}

Temporal evidence is only defined where the aligned memory is itself
usable. We therefore define the per-pixel \emph{support} as the
conjunction of the current raw validity, the aligned temporal validity,
and the warp support,
%
\begin{equation}
S_t
=
M_t^{\raw}
\wedge
\widetilde M_t
\wedge
\mathcal{S}_t ,
\label{eq:support}
\end{equation}
%
and confine the intervention to it. Final disparity is
%
\begin{equation}
\boxed{
d_t^{\fused}
=
\begin{cases}
(1-w_t)d_t^{\raw}+w_t\widetilde m_t, & S_t=1,\\[2pt]
d_t^{\raw}, & S_t=0.
\end{cases}
}
\label{eq:fusion}
\end{equation}

Because $w_t\in[0,1]$, on the support
%
\begin{equation}
\min(d_t^{\raw},\widetilde m_t)
\leq
d_t^{\fused}
\leq
\max(d_t^{\raw},\widetilde m_t),
\end{equation}
%
and off the support the module is the exact identity on raw geometry.

The refiner therefore cannot synthesize arbitrary disparity outside the
interval spanned by its two candidates.

Applying \eqref{eq:support} to the output --- rather than only computing
it --- is a load-bearing part of the contract, not a formality. The
causal pull-warp uses zero padding, so an out-of-bounds sample returns
$\widetilde m_t=0$ \emph{exactly}. An unmasked convex blend at such a
pixel can therefore drive $d_t^{\fused}$ to zero, producing an invalid
disparity where the frozen stereo estimator produced a valid one, and
the bounded recurrence of Sec.~\ref{sec:recurrence} then carries that
zero forward until the next raw re-anchor. Section~\ref{sec:masking}
quantifies the effect: it is invisible on SCARED-C and dominates the
external-domain metric-depth tail.

The exact trainable parameter count is 177,338, and masking changes no
parameter, no threshold, and no training procedure.

Reset output bias is initialized to $-2$, while fusion is initialized
with zero bias, producing conservative initial temporal intervention.


\subsection{Bounded Corrected-State Recurrence}
\label{sec:recurrence}

After refinement,
%
\begin{equation}
m_t=d_t^{\fused}.
\end{equation}

This creates the feedback path
%
\begin{equation}
d_{t-1}^{\fused}
\rightarrow
\widetilde m_t
\rightarrow
d_t^{\fused}
\rightarrow
\widetilde m_{t+1}.
\end{equation}

Continuous propagation can therefore transform previous alignment and
fusion errors into future evidence.

We explicitly bound corrected-state lifetime.
On a sequence boundary the first output is raw.
After a re-anchor, previous raw disparity is used as state.
Otherwise, the previous fused output is propagated.

The canonical policy permits four corrected-state uses before a raw
re-anchor.
Thus $H$ denotes state lifetime rather than a temporal input window.

The method remains one-step causal:
%
\begin{equation}
\{I_{t-1},I_t,m_{t-1}\}
\rightarrow
d_t^{\fused}.
\end{equation}


\subsection{Training Objective}

For target disparity $g_t$, define
%
\begin{align}
e_t^{\raw}
&=
|d_t^{\raw}-g_t|,
\\
e_t^{\mem}
&=
|\widetilde m_t-g_t|,
\\
\Delta_t
&=
e_t^{\mem}-e_t^{\raw}.
\end{align}

Training support is the intersection of GT coverage, current validity,
aligned temporal validity, and warp support.

The disparity reconstruction term is
%
\begin{equation}
\mathcal{L}_{disp}
=
\mu_M
\left[
\operatorname{Huber}_{\delta=1}
(d_t^{\fused}-g_t)
\right].
\end{equation}

Native reset and fusion margins are respectively
%
\begin{equation}
\tau_r^{n}=5\;\mathrm{px},
\qquad
\tau_f^{n}=1\;\mathrm{px},
\end{equation}
%
corresponding at cache scale to
%
\begin{equation}
\tau_r^{c}=0.703125,
\qquad
\tau_f^{c}=0.140625.
\end{equation}

Define
%
\begin{align}
R^+ &= M\cap\{\Delta>\tau_r^c\},\\
R^- &= M\cap\{\Delta<-\tau_r^c\},
\end{align}
%
and
%
\begin{align}
F^+ &= M\cap\{\Delta>\tau_f^c\},\\
F^- &= M\cap\{\Delta<-\tau_f^c\},\\
F^0 &= M\cap\{|\Delta|\leq\tau_f^c\}.
\end{align}

The implemented losses are
%
\begin{align}
\mathcal{L}_{reset}
&=
\mu_{R^+}[r_t]
+
\mu_{R^-}[1-r_t],
\\
\mathcal{L}_{fusion}
&=
\mu_{F^+}[f_t]
+
\mu_{F^-}[1-f_t]
+
0.2\mu_{F^0}[|f_t-0.5|].
\end{align}

The full objective is
%
\begin{equation}
\boxed{
\mathcal{L}
=
\mathcal{L}_{disp}
+
\mathcal{L}_{reset}
+
\mathcal{L}_{fusion}.
}
\end{equation}

No explicit temporal-consistency objective is used.


% ============================================================
% IV. EXPERIMENTAL PROTOCOL
% ============================================================

\section{Experimental Protocol}
\label{sec:experimental}

\subsection{Canonical Training}

The canonical temporal refiner is trained on temporally curated SCARED-C
data derived from SCARED \cite{allan2021scared}.

The local partition is
%
\begin{equation}
\text{train}=\{D1,D3,D6\},
\qquad
\text{development}=D2,
\qquad
\text{held-out}=D7.
\end{equation}

Training uses frozen predictions from:

\begin{itemize}
    \item S$^2$M$^2$-S \cite{min2025s2m2};
    \item RAFT-Stereo \cite{lipson2021raftstereo};
    \item Stereo Anywhere \cite{bartolomei2025stereoanywhere}.
\end{itemize}

CREStereo \cite{li2022crestereo} and Fast-FoundationStereo
\cite{wen2026fastfoundationstereo} are excluded from temporal-head
training.

No backbone identity is given to STEER.


\begin{table}[t]
\centering
\caption{Canonical STEER training configuration.}
\label{tab:train}
\small
\begin{tabular}{ll}
\toprule
Setting & Value \\
\midrule
Train / dev. / held out & D1,D3,D6 / D2 / D7 \\
Training frames & 8,616 \\
Causal pairs & 8,606 \\
Seen backbones & 3 \\
Clip policy & non-overlap, $H=4$ \\
Batch size & 4 \\
Epochs & 12 \\
Optimizer & AdamW \\
Learning rate & $2\times10^{-4}$ \\
Weight decay & $10^{-4}$ \\
Gradient clipping & 1.0 \\
Trainable parameters & 177,338 \\
Model selection & development fused EPE \\
\bottomrule
\end{tabular}
\end{table}


\subsection{External Evaluation}

\paragraph{DRENDS.}
DRENDS contains dynamic robotic-endoscopy sequences with stereo imagery,
calibration, and metric geometry \cite{loza2025drends}.
We report five curated recordings ---
\texttt{Vid10\_Liver\_Med}, \texttt{Vid11\_Liver\_High},
\texttt{Vid12\_Pancreas\_Ext}, \texttt{Vid13\_Pancreas\_Med} and
\texttt{Vid14\_Pancreas\_High} --- evaluated with three stereo
estimators for $20{,}929$ frames in total.
The zero-shot claim applies only to the SCARED-trained canonical
checkpoint: no DRENDS data enters training, checkpoint selection or
threshold calibration.

Disparity reference is derived from metric ToF geometry through stereo
calibration.
It is therefore useful metric evidence but not independent stereo
ground truth, and the ToF reference itself is temporally filtered.

\paragraph{D4D.}
D4D contains deformable surgical stereo sequences and structured-light
observations \cite{docea2026d4d}.
Because the current geometric scale contract remains unresolved, we use
D4D only for no-reference temporal diagnostics.

\paragraph{SERV-CT.}
SERV-CT provides calibrated stereo pairs and CT-derived reference
geometry \cite{edwards2022servct}.
The evaluated pairs are non-consecutive, so temporal refinement is
marked not applicable.


% ============================================================
% UNIFIED EVALUATION
% ============================================================

\subsection{Unified Temporal-Refinement Evaluation}

A central requirement is that evaluation support cannot depend on the
prediction:
%
\begin{equation}
M_{\mathrm{eval}}
=
M_{\mathrm{GT}}
\cap
M_{\mathrm{protocol}}.
\end{equation}

Invalid predictions on valid support are penalized rather than silently
removed.

The framework evaluates five complementary dimensions.

\paragraph{Spatial geometry.}
For disparity we report EPE/MAE, median, RMSE, P90/P95/P99, maximum,
invalid rate, and Bad-$1/3/5$.
For metric depth we additionally report Bad-$2/5/10$ mm, AbsRel, SqRel,
and $\delta_{1,2,3}$.

\paragraph{Temporal fidelity.}
At horizons $k\in\{1,2,4,8\}$ the framework defines grid-based
temporal change error and, when correspondence is available,
motion-compensated $\TEPE_{\mathrm{corr}}$ or warped depth DTCE.
Pure prediction consistency can additionally be measured using OPW.
Estimated-flow correspondences are explicitly marked as
flow-conditioned proxies.

\TODO{
TEMPORAL-METRIC IMPLEMENTATION AUDIT:
the current definitive SCARED-C artifacts contain DTCE_grid and Drift.
Before submission, verify whether TEPE_corr and OPW were actually
executed and persisted elsewhere. If not, describe them as supported
framework metrics rather than reported experimental results.
}

\paragraph{Intervention harm and benefit.}
For
%
\begin{equation}
\delta e=e^{\fused}-e^{\raw},
\end{equation}
%
we measure
%
\begin{align}
\HPlus &= \mathbb{E}[\max(\delta e,0)],\\
\BPlus &= \mathbb{E}[\max(-\delta e,0)],\\
\HUR &= P(\delta e>0),\\
\BUR &= P(\delta e<0).
\end{align}

For each bad-pixel threshold $\tau$ we additionally report
NewBad, RecoveredBad, bad-rate delta, and identity preservation.

\paragraph{Frame-level tails.}
For each frame,
%
\begin{equation}
D_t
=
\frac{1}{|M_t|}
\sum_{x\in M_t}
(e_t^{\fused}(x)-e_t^{\raw}(x)).
\end{equation}

We retain mean, median, P95, P99, worst degradation, and the fraction
of degraded frames.

\paragraph{Selective risk.}
When a risk score is available, the same framework supports
risk--coverage analysis, $\AURC$, oracle/excess $\AURC$, AUROC, AUPRC,
Brier, NLL, and ECE \cite{geifman2019selectivenet}.

Primary aggregation is equal-sequence macro averaging.
Pixel-weighted micro statistics are retained secondarily.


% ============================================================
% V. RESULTS
% ============================================================

\section{Results}
\label{sec:results}


\subsection{Does the Learned Temporal Policy Matter?}

We first isolate the contribution of the learned decision policy on D2.
All variants use the same frozen stereo predictions, temporal alignment,
support contract, and evaluation implementation.
Only the temporal policy is changed.

\begin{table}[t]
\centering
\scriptsize
\caption{
SCARED-C D2 matched temporal-policy closure.
Raw EPE is $0.276202$.
Positive reduction denotes improvement.
}
\label{tab:closure}
\begin{tabular}{lrr}
\toprule
Method & EPE & Reduction \\
\midrule
STEER $H=6$ & \textbf{0.264656} & 4.18\% \\
STEER $H=8$ & 0.264912 & 4.09\% \\
STEER $H=4$ & 0.265028 & 4.05\% \\
STEER $H=2$ & 0.267265 & 3.24\% \\
STEER $H=1$ & 0.269001 & 2.61\% \\
\midrule
Fixed $w=0.5$ & 0.274562 & 0.59\% \\
EMA2 & 0.274562 & 0.59\% \\
Fixed $w=0.3$ & 0.274852 & 0.49\% \\
Fixed $w=0.2$ & 0.275229 & 0.35\% \\
EMA3 & 0.275386 & 0.30\% \\
Fixed $w=0.1$ & 0.275684 & 0.19\% \\
Warped raw previous & 0.276298 & $-0.03\%$ \\
FB confidence weight & 0.282182 & $-2.17\%$ \\
Warped recurrent $w=1$ & 0.286035 & $-3.56\%$ \\
Continuous recurrence & 0.449349 & $-62.69\%$ \\
\midrule
Oracle raw/memory & 0.242223 & 12.30\% \\
\bottomrule
\end{tabular}
\end{table}

The canonical learned policy obtains
%
\begin{equation}
\Delta_{\mathrm{STEER}}
=
0.276202-0.265028
=
0.011174\;\mathrm{px}.
\end{equation}

The best matched fixed fusion obtains only
%
\begin{equation}
\Delta_{\mathrm{fixed}}
=
0.276202-0.274562
=
0.001640\;\mathrm{px}.
\end{equation}

Therefore,
%
\begin{equation}
\boxed{
\frac{
\Delta_{\mathrm{STEER}}
}{
\Delta_{\mathrm{fixed}}
}
\approx
6.8
}
\end{equation}
%
and the learned decision policy recovers almost seven times the
geometric gain of the best fixed temporal blend.

Direct forward--backward confidence is particularly informative:
although $C^{FB}$ is one input cue to STEER, using it directly as the
fusion weight worsens EPE by $2.17\%$.
Thus motion reliability is useful evidence but insufficient as a
decision rule.

Similarly, directly propagating aligned memory or using continuous
corrected recurrence degrades performance.
The improvement is therefore not explained by motion alignment alone,
temporal averaging, or generic smoothing.

Under the implemented one-step update definition, EMA2 is algebraically
equivalent to fixed $w=0.5$; their identical outputs are therefore
expected and should not be interpreted as two independent confirmations.


\subsection{Available Temporal Headroom}

The pixelwise oracle chooses, using ground truth, the better of current
raw disparity and aligned temporal memory.

It reaches
%
\begin{equation}
0.276202
\rightarrow
0.242223,
\end{equation}
%
corresponding to a potential $12.30\%$ reduction.

The oracle gain is
%
\begin{equation}
\Delta_{\mathrm{oracle}}
=
0.033979\;\mathrm{px}.
\end{equation}

Canonical $H=4$ therefore captures
%
\begin{equation}
\boxed{
\frac{
0.011174
}{
0.033979
}
=
32.9\%
}
\end{equation}
%
of the available raw-versus-memory oracle improvement.

This exposes an important asymmetry:
the aligned memory already contains considerable useful information,
but identifying where to use it remains the main unresolved problem.


\subsection{Why $H=4$? Risk--Gain Rather Than Minimum EPE}

The lowest development EPE occurs at $H=6$, not $H=4$:
%
\begin{equation}
0.264656
\quad\text{vs.}\quad
0.265028.
\end{equation}

The absolute difference is only
%
\begin{equation}
0.000372\;\mathrm{px},
\end{equation}
%
approximately $0.14\%$ of the $H=4$ EPE.

However, longer state lifetime increases intervention risk.

\begin{table}[t]
\centering
\scriptsize
\caption{
D2 risk--gain comparison for the strongest bounded horizons.
}
\label{tab:horizon}
\begin{tabular}{lrrrrr}
\toprule
& EPE & Deg. frames & HUR & BUR & $B^+/H^+$ \\
\midrule
$H=4$ & 0.265028 & \textbf{41.46\%} & \textbf{0.4048} & \textbf{0.5074} & \textbf{2.53} \\
$H=6$ & \textbf{0.264656} & 42.18\% & 0.4053 & 0.5065 & 2.30 \\
$H=8$ & 0.264912 & 42.18\% & 0.4068 & 0.5047 & 2.04 \\
\bottomrule
\end{tabular}
\end{table}

Tail behaviour also favours $H=4$:
%
\begin{align}
P95(D_t):
&\quad
0.01715,\;
0.02037,\;
0.02441,
\\
\mathrm{NewBad1}:
&\quad
0.0467\%,\;
0.0676\%,\;
0.1282\%
\end{align}
%
for $H=4,6,8$, respectively.

We therefore do not claim that four recurrent uses minimize
development EPE.
Instead, $H=4$ is the frozen risk--gain operating point:
it sacrifices $0.000372$ px relative to the minimum-EPE $H=6$ policy
while reducing intervention tails.


\subsection{Why Recurrence Must Be Bounded}

Continuous corrected-state recurrence changes mean EPE from
%
\begin{equation}
0.276202
\rightarrow
0.449349,
\end{equation}
%
corresponding to a $62.69\%$ increase in error.

This demonstrates that, for the implemented recurrent state,
continuously writing corrected disparity back into memory accumulates
alignment and fusion errors.

Periodic raw re-anchoring is therefore not a minor implementation
detail but an essential structural component.

This result does not imply that every possible continuous recurrent
architecture must fail.


\subsection{Cross-Backbone Development Behaviour}

Canonical $H=4$ improves all five stereo estimators in mean D2 EPE.

\begin{table}[t]
\centering
\small
\caption{Canonical $H=4$ D2 gain by upstream stereo estimator.}
\label{tab:d2_backbones}
\begin{tabular}{lcc}
\toprule
Backbone & Status & EPE reduction \\
\midrule
Stereo Anywhere & seen & 8.44\% \\
CREStereo & unseen & 3.37\% \\
RAFT-Stereo & seen & 2.87\% \\
Fast-FoundationStereo & unseen & 2.85\% \\
S$^2$M$^2$-S & seen & 2.17\% \\
\bottomrule
\end{tabular}
\end{table}

Seen backbones improve by $4.50\%$ on average, while unseen backbones
improve by $3.11\%$.
All six unseen-backbone sequence groups improve.

The signal is strongest on Stereo Anywhere, but does not originate from
that backbone alone.


\subsection{Held-Out Cross-Backbone Confirmation}

We next evaluate the already frozen $H=4$ policy on held-out D7.

\begin{table*}[t]
\centering
\caption{
Cross-backbone evaluation on SCARED-C.
The same frozen STEER checkpoint is used in every row.
EPE, Bad1, and RMSE use the primary macro-sequence aggregation.
Lower is better.
}
\label{tab:scared_main}
\scriptsize
\setlength{\tabcolsep}{3.7pt}
\begin{tabular}{llccc|ccc}
\toprule
& &
\multicolumn{3}{c|}{D2 development}
&
\multicolumn{3}{c}{D7 held out}
\\
\cmidrule(lr){3-5}
\cmidrule(lr){6-8}
Backbone & Status &
EPE raw$\rightarrow$STEER &
Bad1 raw$\rightarrow$STEER &
RMSE raw$\rightarrow$STEER &
EPE raw$\rightarrow$STEER &
Bad1 raw$\rightarrow$STEER &
RMSE raw$\rightarrow$STEER \\
\midrule
S$^2$M$^2$-S
& seen
& .2704$\rightarrow$\textbf{.2645}
& .01394$\rightarrow$\textbf{.01291}
& .3892$\rightarrow$\textbf{.3806}
& .5176$\rightarrow$\textbf{.5010}
& .06440$\rightarrow$\textbf{.06076}
& 1.4515$\rightarrow$\textbf{1.3991} \\
RAFT-Stereo
& seen
& .2691$\rightarrow$\textbf{.2614}
& .01603$\rightarrow$\textbf{.01492}
& .4009$\rightarrow$\textbf{.3896}
& .4671$\rightarrow$\textbf{.4397}
& .06085$\rightarrow$\textbf{.05845}
& 1.3324$\rightarrow$\textbf{1.2290} \\
Stereo Anywhere
& seen
& .3022$\rightarrow$\textbf{.2767}
& .02183$\rightarrow$\textbf{.01754}
& .7208$\rightarrow$\textbf{.4785}
& .4774$\rightarrow$\textbf{.4449}
& .06227$\rightarrow$\textbf{.05828}
& 1.2869$\rightarrow$\textbf{1.1895} \\
CREStereo
& unseen
& .2668$\rightarrow$\textbf{.2578}
& .01568$\rightarrow$\textbf{.01439}
& .4601$\rightarrow$\textbf{.4299}
& .5165$\rightarrow$\textbf{.4957}
& .06354$\rightarrow$\textbf{.06108}
& 1.3757$\rightarrow$\textbf{1.3120} \\
Fast-FoundationStereo
& unseen
& .2725$\rightarrow$\textbf{.2647}
& .01604$\rightarrow$\textbf{.01477}
& .4022$\rightarrow$\textbf{.3905}
& .4819$\rightarrow$\textbf{.4668}
& .05919$\rightarrow$\textbf{.05684}
& 1.3535$\rightarrow$\textbf{1.3076} \\
\bottomrule
\end{tabular}
\end{table*}

All five backbones improve on D7, including both architectures unseen
during temporal-head training.

The mean backbone EPE decreases from
%
\begin{equation}
0.4921
\rightarrow
0.4696,
\end{equation}
%
or $4.57\%$.

The improvement is not restricted to the mean.
Bad1 and disparity RMSE decrease for every backbone on both D2 and D7.
P99 disparity error likewise decreases in all ten backbone--split cases,
while maximum error decreases in nine of ten.
The strongest tail reduction occurs for Stereo Anywhere on D2, where
maximum disparity error decreases from $33.27$ to $14.88$ px.
InvalidRate remains zero for both raw and refined outputs in all ten
SCARED-C backbone--split evaluations. This is a property of SCARED-C
rather than of the module in general: Section~\ref{sec:masking} shows
that on external domains an unmasked output does introduce invalid
predictions, which is precisely what the support contract of
\eqref{eq:fusion} prevents.

Thus the EPE gain is not obtained by trading mean accuracy for a worse
aggregate bad-pixel rate or a broader SCARED-C error tail.
This does not imply that no individual new bad pixels are introduced;
those events are measured explicitly through NewBad and recovered-error
statistics below.

This supports transfer of one common learned intervention policy across
heterogeneous stereo estimators within the SCARED-C domain.

The joint unseen-backbone / unseen-domain cell is reported separately in
Section~\ref{sec:drends}.


\subsection{Refinement as an Intervention}

The canonical model does not merely reduce final EPE; its beneficial
updates consistently outweigh the new threshold crossings that it
introduces.
Table~\ref{tab:intervention_quality} summarizes two complementary
quantities:
the beneficial-to-harmful error-magnitude ratio $B^+/H^+$ and the ratio
between recovered and newly introduced Bad1 pixels.

\begin{table}[t]
\centering
\scriptsize
\caption{
Canonical STEER intervention quality.
All ratios exceed one; higher is better.
``Rec./New'' is RecoveredBad1/NewBad1.
}
\label{tab:intervention_quality}
\begin{tabular}{lrrrr}
\toprule
& \multicolumn{2}{c}{D2} & \multicolumn{2}{c}{D7} \\
\cmidrule(lr){2-3}\cmidrule(lr){4-5}
Backbone & $B^+/H^+$ & Rec./New & $B^+/H^+$ & Rec./New \\
\midrule
S$^2$M$^2$-S          & 1.76 & 3.1$\times$  & 2.10 & 6.0$\times$ \\
RAFT-Stereo           & 2.05 & 2.9$\times$  & 2.85 & 3.2$\times$ \\
Stereo Anywhere       & 4.59 & 10.4$\times$ & 3.30 & 5.9$\times$ \\
CREStereo             & 2.23 & 4.3$\times$  & 2.40 & 4.6$\times$ \\
Fast-FoundationStereo & 2.11 & 4.0$\times$  & 2.10 & 4.8$\times$ \\
\bottomrule
\end{tabular}
\end{table}

Across both splits and all five backbones, STEER recovers between
$2.9\times$ and $10.4\times$ as many Bad1 pixels as it newly introduces.
The corresponding identity-preservation score exceeds $99.87\%$ in
every case.
This directly supports the refinement-centric interpretation:
the method changes only a small subset of threshold status while
recovering substantially more errors than it creates.



\subsection{Grid-Based Multi-Horizon Temporal Diagnostic}

The definitive SCARED-C artifacts also contain grid-based temporal
change error at horizons $k\in\{1,2,4,8\}$.
Because this quantity is evaluated at fixed image-grid locations, it is
\emph{not} motion compensated and should not be interpreted as temporal
fidelity under scene or camera motion.

\begin{table*}[t]
\centering
\scriptsize
\caption{
Relative change in grid-based DTCE MAE after refinement.
Negative values denote lower grid-based temporal-change error.
This is a diagnostic, not a motion-compensated fidelity metric.
}
\label{tab:dtce_grid}
\setlength{\tabcolsep}{4.4pt}
\begin{tabular}{llrrrr}
\toprule
Split & Backbone & $k=1$ & $k=2$ & $k=4$ & $k=8$ \\
\midrule
D2 & CREStereo              & -2.9\%  & +3.9\%  & +0.7\%  & +2.4\% \\
D2 & Fast-FoundationStereo  & +2.6\%  & +8.2\%  & +3.0\%  & +4.5\% \\
D2 & RAFT-Stereo            & -0.4\%  & +6.5\%  & +2.6\%  & +4.3\% \\
D2 & S$^2$M$^2$-S           & +1.9\%  & +7.4\%  & +2.7\%  & +4.1\% \\
D2 & Stereo Anywhere        & -31.6\% & -22.1\% & -20.1\% & -14.8\% \\
\midrule
D7 & CREStereo              & -15.9\% & -7.2\%  & -4.9\%  & -3.6\% \\
D7 & Fast-FoundationStereo  & -9.2\%  & -2.6\%  & -2.5\%  & -2.1\% \\
D7 & RAFT-Stereo            & -18.7\% & -8.9\%  & -5.6\%  & -4.6\% \\
D7 & S$^2$M$^2$-S           & -11.3\% & -5.7\%  & -4.9\%  & -3.2\% \\
D7 & Stereo Anywhere        & -26.2\% & -15.9\% & -12.5\% & -8.6\% \\
\bottomrule
\end{tabular}
\end{table*}

On held-out D7, grid-based DTCE decreases for all five backbones at all
four horizons ($20/20$ backbone--horizon cases).
On D2 the behavior is mixed: several backbones improve geometrically in
EPE, Bad1, and RMSE while grid-based DTCE increases at some horizons.

This contrast is useful rather than contradictory.
It demonstrates directly that improved frame-wise geometry and reduced
grid-based temporal change are not interchangeable objectives.
In particular, a motion-uncompensated grid metric can penalize changes
that reflect legitimate scene or camera motion.

\subsection{External Zero-Shot Geometry on DRENDS}
\label{sec:drends}

We evaluate the frozen canonical checkpoint on five DRENDS recordings
with three stereo estimators, two of which were excluded from
temporal-head training, for a total of $20{,}929$ frames. No component
is adapted, retrained, or tuned for this domain.

\begin{table*}[t]
\centering
\scriptsize
\caption{
Zero-shot DRENDS transfer, relative change after refinement
(negative is better).
The reference disparity is derived from temporally filtered ToF
geometry and is metric evidence, not independent stereo ground truth.
The two rightmost estimators were excluded from temporal-head training,
so their columns are the joint unseen-backbone / unseen-domain setting.
}
\label{tab:drends}
\setlength{\tabcolsep}{5pt}
\begin{tabular}{lcccccc}
\toprule
& \multicolumn{2}{c}{RAFT-Stereo (seen)}
& \multicolumn{2}{c}{CREStereo (unseen)}
& \multicolumn{2}{c}{Fast-FoundationStereo (unseen)} \\
\cmidrule(lr){2-3}\cmidrule(lr){4-5}\cmidrule(lr){6-7}
Recording & EPE & Bad3 & EPE & Bad3 & EPE & Bad3 \\
\midrule
Vid10\_Liver\_Med        & -5.7\% & -11.3\% & -5.2\% & -8.7\%  & -5.0\% & -7.3\% \\
Vid11\_Liver\_High       & -3.4\% & -21.9\% & -3.7\% & -22.8\% & -3.2\% & -20.4\% \\
Vid12\_Pancreas\_Ext     & -8.4\% & -22.7\% & -6.3\% & -18.5\% & -5.3\% & -10.3\% \\
Vid13\_Pancreas\_Med     & -4.3\% & -14.2\% & -5.2\% & -10.4\% & -4.3\% & -10.9\% \\
Vid14\_Pancreas\_High    & -4.7\% & -27.1\% & -5.6\% & -26.5\% & -4.7\% & -27.4\% \\
\bottomrule
\end{tabular}
\end{table*}

Across the 15 recording--backbone combinations, disparity EPE, Bad3,
metric depth MAE and metric depth RMSE all improve in \emph{every} cell
($60/60$), with mean changes of $-5.00\%$, $-17.35\%$, $-3.68\%$ and
$-3.21\%$ respectively.

Two observations matter more than the aggregate.
First, the per-recording pattern is set by the \emph{recording} and not
by the estimator: on Vid10 the EPE change is $-5.7\%$, $-5.2\%$ and
$-5.0\%$ for the seen and the two unseen backbones respectively, and the
same alignment holds on every other recording. The residual variance is
therefore a property of the data, which is stronger evidence for
backbone independence than any pooled mean.
Second, both unseen estimators improve on all five recordings, so the
joint unseen-backbone / unseen-domain cell is positive rather than
merely unevaluated.

The transfer matrix is consequently complete:
seen-backbone / seen-domain and unseen-backbone / seen-domain on
SCARED-C (Table~\ref{tab:scared_main}), and seen-backbone / OOD-domain
and unseen-backbone / OOD-domain on DRENDS
(Table~\ref{tab:drends}).
This establishes zero-shot geometric transfer in disparity space under
a ToF-derived metric reference; it does not establish universal
cross-domain robustness, and DRENDS remains a non-independent reference.


\subsection{The Support Contract Is Load-Bearing}
\label{sec:masking}

The results above use the support-confined output of
\eqref{eq:fusion}. Because the ablation is a single line, it isolates
the contract cleanly: the checkpoint, the head, the evidence, the
policy and the reset schedule are identical, and only the masking of
the output differs.

\begin{table}[t]
\centering
\scriptsize
\caption{
Effect of applying the support mask of \eqref{eq:support} to the module
output, over five DRENDS recordings and three backbones.
``Regresses'' counts recording--backbone cells that get worse.
}
\label{tab:masking}
\begin{tabular}{lrr}
\toprule
Quantity & Unmasked output & Support-confined \\
\midrule
Invalid refined pixels & up to $1.3\times10^{-5}$ & \textbf{0} \\
Depth RMSE & regresses $14/14$, $+110.4\%$ & \textbf{improves $15/15$, $-3.2\%$} \\
Depth MAE  & regresses $10/14$, $+4.2\%$   & \textbf{improves $15/15$, $-3.7\%$} \\
Disparity EPE  & $-4.30\%$  & \textbf{$-5.00\%$} \\
Disparity Bad3 & $-14.06\%$ & \textbf{$-17.35\%$} \\
\bottomrule
\end{tabular}
\end{table}

Without the mask, the causal pull-warp's zero padding lets an
out-of-bounds sample enter the convex blend as $\widetilde m_t=0$, the
fused disparity collapses to zero at those pixels, and the bounded
recurrence propagates the collapse until the next raw re-anchor. The
offending pixels are rare and bursty --- $300$ of $38{,}828{,}160$ on
Vid10\_Liver\_Med and $265$ of $38{,}439{,}360$ on Vid11\_Liver\_High,
concentrated in runs of at most $H$ consecutive frames --- and the
aligned memory of every one of them is exactly $0.0$. Bounded
recurrence therefore contains the damage by construction but does not
prevent its injection.

Because an invalid prediction on valid support is penalized rather than
discarded, this rare event dominates the metric-depth tail: it is the
entire source of the depth-RMSE reversal previously attributed to the
external domain. Confining the output removes it, and improves
disparity as well, since the zero no longer enters memory.

On SCARED-C the same ablation is almost inert: D2 pooled EPE gain is
unchanged at $4.05\%$, while D7 falls slightly from $4.64\%$ to
$4.57\%$, with the largest single change on RAFT-Stereo
($6.10\%\rightarrow5.87\%$). The support contract is therefore not free
in-domain; it costs a small amount of in-domain gain and buys
well-posed behaviour out of domain. We report the confined variant
because a module that declares a support must respect it.


\subsection{No-Reference Temporal Behaviour on D4D}

Because the current D4D metric-geometry contract remains unresolved,
we report only prediction-space temporal inconsistency.

\begin{table}[t]
\centering
\caption{
D4D motion-compensated prediction inconsistency.
This is not a geometry metric.
}
\label{tab:d4d}
\begin{tabular}{lrrr}
\toprule
Backbone & Raw & STEER & Reduction \\
\midrule
RAFT-Stereo
& .5094 & \textbf{.2450} & 51.9\% \\
Stereo Anywhere
& 6.9271 & \textbf{3.1516} & 54.5\% \\
\bottomrule
\end{tabular}
\end{table}

Prediction inconsistency decreases substantially.
However, a degenerate smoother can also obtain low temporal
inconsistency while preserving stale geometry.
We therefore make no D4D geometric claim.


\subsection{Larger Training Does Not Improve the Operating Point}

A larger augmented-training run retains the same architecture.

\begin{table}[t]
\centering
\scriptsize
\caption{
Held-out D7 canonical versus augmented training.
}
\label{tab:augmented}
% NOTE: Degraded-frame fraction must come from PositiveFraction.
% Do not substitute CatastrophicFraction until the metric wiring audit
% is complete.
\begin{tabular}{lrr}
\toprule
Metric & Canonical & Augmented \\
\midrule
Mean EPE & \textbf{.4693} & .4706 \\
$\HUR$ & \textbf{.3965} & .4457 \\
$\HPlus$ & \textbf{.0148} & .0235 \\
$\BPlus$ & .0376 & \textbf{.0450} \\
$\BUR$ & .4951 & \textbf{.5060} \\
Degraded-frame fraction & \textbf{.2642} & .3341 \\
Worst frame degradation & \textbf{.1063} & .3958 \\
\bottomrule
\end{tabular}
\end{table}

The augmented model performs more beneficial updates, but also
substantially more harmful ones.
It improves raw stereo, yet does not outperform the canonical operating
point overall.

This suggests that increasing training scale can shift intervention
aggressiveness without improving held-out risk--gain balance.


\subsection{External Temporal Stabilizer Context}

BiDAStabilizer is the closest plug-in stereo stabilizer in our current
comparison set \cite{jing2024bida}.
Unlike STEER, its full sequence formulation is non-causal because it
uses future-frame/bidirectional temporal propagation.

On the full SCARED-C D2 sequence set, our independent reproduction
changes pixel-weighted EPE from
%
\begin{equation}
0.302862
\rightarrow
0.302109.
\end{equation}

Because the BiDA and STEER evaluation paths currently use
independently validated but not pixel-identical support contracts, we
do not directly subtract their absolute values.

\TODO{
SUBMISSION-CRITICAL COMPARISON:
generate a pixel-identical/common-support RAFT-Stereo table for
Raw vs BiDAStabilizer vs STEER, including EPE, Bad1/Bad3, HUR,
HPlus, BPlus, NewBad, temporal fidelity, and worst-frame degradation.
Keep the causal/non-causal distinction explicit.
Until this table exists, the independent BiDA number above is context
only and must not be used for a direct superiority claim.
}

NVDS is used only as a relative-depth temporal-stability reference
\cite{wang2023nvds}.


% ============================================================
% VI. ANALYSIS
% ============================================================

\section{Analysis}
\label{sec:analysis}


\subsection{What Actually Makes STEER Work?}

The controlled closure supports a specific interpretation.

First, geometric alignment is necessary but not sufficient.
Aligned temporal memory contains useful information, as shown by the
$12.30\%$ oracle ceiling.

Second, using this memory indiscriminately fails.
Warped raw memory produces no meaningful gain; direct FB-confidence
weighting and full recurrent replacement worsen EPE.

Third, the learned decision policy matters.
Under identical temporal evidence and support, canonical STEER obtains
almost seven times the gain of the best fixed blend.

Fourth, recurrent state must remain bounded.
Continuous corrected-state recurrence accumulates error and eventually
collapses.

The empirical mechanism is therefore:

\begin{equation}
\boxed{
\begin{gathered}
\text{motion alignment}
+
\text{learned graded intervention}
\\[-0.5mm]
+
\text{bounded recurrent state}
\end{gathered}
}
\end{equation}

rather than generic temporal smoothing.


\subsection{Why Can a 177k-Parameter Head Generalize?}

The temporal head does not solve stereo matching.

It receives two already meaningful geometric candidates:
%
\begin{equation}
d_t^{\raw}
\qquad\text{and}\qquad
\widetilde m_t,
\end{equation}
%
along with evidence describing their disagreement, motion reliability,
validity, matching consistency, and appearance.

The learning problem is therefore closer to
%
\begin{equation}
P
\left(
\text{temporal evidence is useful}
\mid
\text{external reliability cues}
\right)
\end{equation}
%
than to predicting disparity from RGB.

Convex fusion additionally restricts the action space.
The model must decide how strongly to move between two candidate
geometries rather than synthesize arbitrary depth.

This constrained decision problem provides a plausible explanation for
why one small head transfers to stereo architectures that were not used
during temporal-head training.


\subsection{Oracle Gap as the Next Problem}

STEER $H=4$ captures only $32.9\%$ of the raw-versus-memory oracle gain.

Thus the principal remaining opportunity is not necessarily richer
temporal memory.
Useful evidence already exists.

The open problem is better discrimination between helpful and harmful
temporal intervention.

This observation motivates future selective/risk-aware formulations
that use the internal reset/fusion signal to estimate intervention
reliability.


\subsection{Relationship to Uncertainty}

STEER does not currently produce a calibrated predictive uncertainty.

However, several inputs and internal outputs are reliability-related:
forward--backward consistency, warp validity, raw--memory disagreement,
matching evidence, and the learned effective temporal weight
%
\begin{equation}
w_t=r_tf_t.
\end{equation}

The D2 closure shows that these cues are more useful jointly than when
FB confidence is used directly as a heuristic weight.

We ran this analysis. On the prediction-independent protocol support we
recover the effective weight exactly from the convex form,
$w_t=(d_t^{\fused}-d_t^{\raw})/(\widetilde m_t-d_t^{\raw})$, and score it
as a predictor of harmful intervention $\delta e>0$.

\begin{table}[t]
\centering
\scriptsize
\caption{
$w_t$ and the two other freely available cues as risk scores for
harmful intervention. AUROC of $0.5$ is chance.
}
\label{tab:risk}
\begin{tabular}{lcc}
\toprule
Score & AUROC on D2 & AUROC on D7 \\
\midrule
Effective weight $w_t$                      & $0.366$--$0.382$ & $0.496$--$0.506$ \\
Disagreement $|\widetilde m_t-d_t^{\raw}|$  & $0.499$--$0.509$ & $0.450$--$0.459$ \\
Update $|d_t^{\fused}-d_t^{\raw}|$          & $0.392$--$0.406$ & $0.474$--$0.489$ \\
\bottomrule
\end{tabular}
\end{table}

On held-out D7, $w_t$ is at chance. The reliability diagram explains
why: for RAFT-Stereo the harm rate \emph{falls} with $w_t$ on D2
($0.550\rightarrow0.324$) but \emph{rises} on D7
($0.260\rightarrow0.450$). The ordering inverts between development and
held-out data, so $w_t$ is not a weak risk score but a
\emph{sign-unstable} one, and it does not transfer. Risk--coverage
confirms this: selecting by $w_t$ captures only $3.5\%$ (D2) and $6.1\%$
(D7) of the achievable oracle selective gain.

$w_t$ is therefore an intervention magnitude and not a calibrated
uncertainty score, and closing the oracle gap will require evidence the
head does not currently receive rather than a better read-out of what
it already computes.


\subsection{Evaluation Matters}

The proposed evaluation framework exposes several behaviours that mean
EPE alone would hide.

On DRENDS, mean and percentile geometry improve while RMSE worsens.

The augmented model increases both beneficial and harmful intervention.

Longer horizons slightly improve mean D2 EPE while degrading
$B^+/H^+$, NewBad1, and frame tails.

These examples motivate evaluating temporal refinement as an
\emph{intervention} rather than only as a final disparity estimator.


One reported quantity required an audit. FrameDegradation's
CatastrophicFraction was numerically identical to PositiveFraction on
every inspected row. The cause is definitional rather than a wiring
fault: the two are $\Pr(D_t>\tau_{\mathrm{cat}})$ and $\Pr(D_t>0)$, and
the catastrophic threshold defaults to $\tau_{\mathrm{cat}}=0$, which
the evaluators never override. At that setting the two statistics
coincide by construction. We therefore report only the degraded-frame
fraction, explicitly defined as $\Pr(D_t>0)$, and omit
CatastrophicFraction rather than present a duplicate column under a
name that implies a severity threshold it does not carry.


% ============================================================
% VII. DISCUSSION AND LIMITATIONS
% ============================================================

\section{Discussion and Limitations}

\paragraph{Development versus confirmation.}

The matched baseline and horizon closure is performed on D2 development
data.
It supports mechanism analysis and the pre-frozen $H=4$ operating
policy, but does not itself constitute held-out confirmation.

The canonical $H=4$ checkpoint is separately evaluated on D7 without
changing the policy.


\paragraph{Cross-backbone transfer.}

STEER improves five SCARED-C stereo estimators, including two excluded
from temporal-head training.
This supports same-domain unseen-backbone transfer.

It does not establish universal backbone agnosticism.


\paragraph{Joint unseen-backbone and OOD transfer.}

All four cells of the transfer matrix are now populated:
seen-backbone / seen-domain and unseen-backbone / seen-domain on
SCARED-C, and seen-backbone / OOD-domain and unseen-backbone /
OOD-domain on DRENDS (Section~\ref{sec:drends}).

Two caveats bound this. The DRENDS reference is temporally filtered ToF
geometry rendered through stereo calibration, so it is metric evidence
rather than independent stereo ground truth; and a single external
domain does not establish universal cross-domain robustness. What the
matrix supports is that the two shifts compose without the intervention
policy degrading, not that the policy is domain-agnostic.


\paragraph{Internal disocclusion.}

Pull-warp support detects out-of-frame sampling but cannot fully model
internal disocclusion.
Newly visible anatomy may inherit stale temporal geometry despite
plausible motion consistency.

This is a plausible source of residual tail failures.


\paragraph{No-reference temporal metrics.}

Prediction smoothness can be gamed by copying or strongly smoothing
historical geometry.
No-reference temporal scores are therefore diagnostics rather than
evidence of geometric correctness.

\TODO{
SUBMISSION-CRITICAL D4D CONTROL:
add copy-previous and warped-previous degenerate temporal baselines.
Without them, no-reference consistency remains diagnostic only.
}


\paragraph{Statistical uncertainty.}

The definitive held-out table currently reports the canonical seed-0
checkpoint.
Additional canonical seeds and sequence-aware confidence intervals are
required for final statistical uncertainty.

\TODO{
SUBMISSION-CRITICAL STATISTICS:
train canonical seeds 1 and 2 with the identical frozen recipe and
checkpoint-selection rule.
Report mean$\pm$std and paired sequence-aware bootstrap confidence
intervals; do not use pixel-wise independence assumptions.
}


\paragraph{Runtime.}

The learned head is compact, but the end-to-end system includes frozen
stereo plus forward and reverse SEA-RAFT.
On an NVIDIA H100 PCIe at the $144\times180$ evidence grid, over 200
timed iterations after 30 warmup iterations, the median per-frame cost
is $13.22$\,ms for SEA-RAFT forward, $13.19$\,ms for SEA-RAFT reverse,
$1.04$\,ms for the alignment, $5.12$\,ms for the 142-channel evidence
including the frozen ResNet-18 features, and $1.03$\,ms for the fusion
head, giving a total temporal overhead of $33.60$\,ms at $124.2$\,MB
peak VRAM.

Added to the frozen stereo latency this yields $11.0$\,FPS with
Fast-FoundationStereo down to $1.4$\,FPS with Stereo Anywhere, so
\emph{we make no real-time claim}.
Two points follow. The $177$k-parameter head accounts for $3\%$ of the
overhead, so ``compact head'' is true but not the operative cost;
$79\%$ is the two frozen flow passes. And the reverse pass costs
$13.19$\,ms while feeding only the forward--backward confidence cue,
whose direct use as a fusion weight \emph{worsens} EPE by $2.17\%$
(Table~\ref{tab:closure}); removing it would cut the temporal overhead
by $39\%$ and is the most promising efficiency ablation.


\paragraph{Downstream robotic relevance.}

Current experiments evaluate perception geometry and temporal behaviour,
not downstream robotic control.
A 3D surface-stability or reconstruction-consistency experiment would
more directly connect temporal refinement to robotic perception.

\TODO{
OPTIONAL HIGH-VALUE ROBOTICS RESULT:
add a compact 3D surface/point-cloud temporal-stability or reconstruction
consistency experiment if time permits.
}


\paragraph{Clinical interpretation.}

The experiments do not demonstrate clinical safety, autonomous surgical
performance, or patient benefit.


% ============================================================
% VIII. CONCLUSION
% ============================================================

\section{Conclusion}

We introduced \steer, a causal plug-and-play temporal refiner for frozen
surgical stereo.

STEER does not relearn stereo.
It aligns historical geometry using frozen motion estimation and learns
a compact intervention policy over a common external evidence space.

A controlled development closure shows that its gain cannot be explained
by simple temporal averaging:
canonical $H=4$ reduces EPE by $4.05\%$, almost seven times the gain of
the strongest matched fixed/EMA baseline.
Direct flow-confidence weighting and indiscriminate memory propagation
instead degrade performance.
The raw-versus-memory oracle exposes a $12.30\%$ improvement ceiling,
showing that useful temporal information remains substantially
underexploited.

Bounded recurrence is equally important:
continuous corrected-state propagation increases EPE by $62.69\%$.
The canonical $H=4$ policy sacrifices only $0.000372$ px relative to
the minimum-EPE $H=6$ development setting while achieving better
risk--gain and frame-tail behaviour.

The frozen $H=4$ checkpoint subsequently improves all five evaluated
stereo estimators on held-out SCARED-C D7, including two unseen during
temporal-head training.
Across D2 and D7, EPE, Bad1, and disparity RMSE improve for every
backbone; recovered Bad1 pixels outnumber newly introduced Bad1 pixels
by $2.9\times$--$10.4\times$.
Grid-based temporal diagnostics further show why geometry and temporal
change must be reported separately: D7 improves at all $20/20$
backbone--horizon combinations, while D2 remains mixed despite consistent
geometric gains.
Preliminary zero-shot DRENDS evaluation further shows OOD-domain
geometric transfer with a training-seen backbone, although a
domain-specific depth-RMSE tail remains.

Finally, the unified temporal-refinement evaluation framework makes
explicit that a useful temporal system must be judged not only by final
accuracy, but also by how much it helps, what it damages, how its errors
evolve over time, and whether rare failures become worse.

The resulting picture is simple:
STEER works not through brute-force model scale, but through
\emph{geometric alignment, learned selective intervention, and bounded
temporal memory}.


% ============================================================
% REFERENCES
% ============================================================

\balance

\begin{thebibliography}{99}

\bibitem{li2023codd}
Z.~Li, W.~Ye, D.~Wang, F.~X. Creighton, R.~H. Taylor,
G.~Venkatesh, and M.~Unberath,
``Temporally consistent online depth estimation in dynamic scenes,''
in \emph{Proc. IEEE/CVF Winter Conf. Applications of Computer Vision
(WACV)}, 2023, pp.~3018--3027.

\bibitem{karaev2023dynamicstereo}
N.~Karaev, I.~Rocco, B.~A.~T. Graham, N.~Neverova, A.~Vedaldi,
and C.~Rupprecht,
``DynamicStereo: Consistent dynamic depth from stereo videos,''
in \emph{Proc. IEEE/CVF Conf. Computer Vision and Pattern Recognition
(CVPR)}, 2023, pp.~13229--13239.

\bibitem{jing2024bida}
J.~Jing, Y.~Mao, A.~Qiu, and K.~Mikolajczyk,
``Match Stereo Videos via Bidirectional Alignment,''
in \emph{Proc. European Conf. Computer Vision (ECCV)}, 2024.

\bibitem{wang2025ppm}
Y.~Wang, J.~Hu, Q.~Dong, Y.~Zhang, Y.~Fu, T.~L.~Lam, and D.~Wu,
``PPMStereo: Pick-and-play memory construction for consistent dynamic
stereo matching,''
in \emph{Advances in Neural Information Processing Systems},
2025.

\bibitem{wang2024searaft}
Y.~Wang, L.~Lipson, and J.~Deng,
``SEA-RAFT: Simple, efficient, accurate RAFT for optical flow,''
in \emph{Proc. European Conf. Computer Vision (ECCV)}, 2024.

\bibitem{he2016resnet}
K.~He, X.~Zhang, S.~Ren, and J.~Sun,
``Deep residual learning for image recognition,''
in \emph{Proc. IEEE Conf. Computer Vision and Pattern Recognition
(CVPR)}, 2016, pp.~770--778.

\bibitem{lipson2021raftstereo}
L.~Lipson, Z.~Teed, and J.~Deng,
``RAFT-Stereo: Multilevel recurrent field transforms for stereo
matching,''
in \emph{Proc. Int. Conf. 3D Vision (3DV)}, 2021, pp.~218--227.

\bibitem{li2022crestereo}
J.~Li, P.~Wang, P.~Xiong, T.~Cai, Z.~Yan, L.~Yang,
J.~Liu, H.~Fan, and S.~Liu,
``Practical stereo matching via cascaded recurrent network with
adaptive correlation,''
in \emph{Proc. IEEE/CVF Conf. Computer Vision and Pattern Recognition
(CVPR)}, 2022, pp.~16263--16272.

\bibitem{bartolomei2025stereoanywhere}
L.~Bartolomei, F.~Tosi, M.~Poggi, and S.~Mattoccia,
``Stereo Anywhere: Robust zero-shot deep stereo matching even where
either stereo or mono fail,''
in \emph{Proc. IEEE/CVF Conf. Computer Vision and Pattern Recognition
(CVPR)}, 2025.

\bibitem{min2025s2m2}
J.~Min, Y.~Jeon, J.~Kim, and M.~Choi,
``S$^2$M$^2$: Scalable stereo matching model for reliable depth
estimation,''
in \emph{Proc. IEEE/CVF Int. Conf. Computer Vision (ICCV)}, 2025.

\bibitem{wen2025foundationstereo}
B.~Wen, M.~Trepte, J.~Aribido, J.~Kautz, O.~Gallo,
and S.~Birchfield,
``FoundationStereo: Zero-shot stereo matching,''
in \emph{Proc. IEEE/CVF Conf. Computer Vision and Pattern Recognition
(CVPR)}, 2025.

\bibitem{wen2026fastfoundationstereo}
B.~Wen, S.~Dewan, and S.~Birchfield,
``Fast-FoundationStereo: Real-time zero-shot stereo matching,''
in \emph{Proc. IEEE/CVF Conf. Computer Vision and Pattern Recognition
(CVPR)}, 2026.

\bibitem{wang2023nvds}
Y.~Wang, M.~Shi, J.~Li, C.~Hong, Z.~Huang, J.~Peng,
Z.~Cao, J.~Zhang, K.~Xian, and G.~Lin,
``NVDS+: Towards efficient and versatile neural stabilizer for video
depth estimation,''
2023, arXiv:2307.08695.

\bibitem{geifman2019selectivenet}
Y.~Geifman and R.~El-Yaniv,
``SelectiveNet: A deep neural network with an integrated reject
option,''
in \emph{Proc. International Conference on Machine Learning (ICML)},
2019.

\bibitem{allan2021scared}
M.~Allan, J.~McLeod, C.~Wang, J.-C.~Rosenthal, Z.~Hu,
N.~Gard, P.~Eisert, K.~X.~Fu, T.~Zeffiro, W.~Xia,
Z.~Zhu, H.~Luo, F.~Jia, X.~Zhang, X.~Li, L.~Sharan,
T.~Kurmann, S.~Schmid, R.~Sznitman, D.~Psychogyios,
M.~Azizian, D.~Stoyanov, L.~Maier-Hein, and S.~Speidel,
``Stereo correspondence and reconstruction of endoscopic data
challenge,''
2021, arXiv:2101.01133.

\bibitem{edwards2022servct}
P.~J. Edwards, D.~Psychogyios, S.~Speidel, L.~Maier-Hein,
and D.~Stoyanov,
``SERV-CT: A disparity dataset from cone-beam CT for validation of
endoscopic 3D reconstruction,''
\emph{Medical Image Analysis},
vol.~76, p.~102302, 2022.

\bibitem{long2021edssr}
Y.~Long, Z.~Li, C.~H.~Yee, C.~F.~Ng, R.~H.~Taylor,
M.~Unberath, and Q.~Dou,
``E-DSSR: Efficient dynamic surgical scene reconstruction with
transformer-based stereoscopic depth perception,''
2021, arXiv:2107.00229.

\bibitem{docea2026d4d}
R.~Docea, R.~Younis, Y.~Long, M.~Fleury, J.~Xu, C.~Li,
A.~Schulze, A.~Wierick, J.~Bender, M.~Pfeiffer, Q.~Dou,
M.~Wagner, and S.~Speidel,
``The Dresden Dataset for 4D reconstruction of non-rigid abdominal
surgical scenes,''
2026, arXiv:2603.02985.

\bibitem{loza2025drends}
G.~Loza, M.~Magro, and B.~Calm\'e,
``DRENDS: A dataset for depth in robotic endoscopy with dynamic
scenarios,''
Zenodo Dataset, 2025,
doi: 10.5281/zenodo.17598453.

\end{thebibliography}

\end{document}